\documentclass[11pt,leqno]{amsart}
\usepackage{hyperref}
\usepackage{epsfig}
\usepackage{enumitem}
\usepackage{amsmath, amssymb}
\usepackage{amscd}
\usepackage[all,cmtip,matrix,arrow]{xy}
\usepackage{mathrsfs}
\usepackage{graphicx}
\usepackage{tikz-cd}
\usepackage{mathabx}
\usepackage{array}
\usepackage{longtable}
\xyoption{rotate}
\xyoption{pdf}
\setcounter{section}{-1}
%\usepackage[all,cmtip]{xy}
%\usepackage{gclc}
%\usepackage[russian]{babel}
\usepackage{xparse}

\sloppy

\newcommand{\Implies}[2]{$\text{\ref{#1}}\implies\text{\ref{#2}}$}
\newcommand{\Iff}[2]{$\text{\ref{#1}}\iff\text{\ref{#2}}$}
\newcommand{\Ifff}[3]{$\text{\ref{#1}}\iff\text{\ref{#2}}\iff\text{\ref{#3}}$}

\makeatletter
\DeclareRobustCommand*{\bfseries}{%
  \not@math@alphabet\bfseries\mathbf
  \fontseries\bfdefault\selectfont
  \boldmath
}
\makeatother

\setlength{\textheight}{8.5in} \setlength{\textwidth}{6in}
\setlength{\topmargin}{0in} \setlength{\oddsidemargin}{0.25in}
\setlength{\evensidemargin}{0.25in}

\mathsurround=2pt
\AtBeginEnvironment{tikzcd}{\setlength{\mathsurround}{0pt}}
\renewcommand{\baselinestretch}{1.3}

\sloppy

\newtheorem{theo}{Theorem}[section]
\newtheorem{lemma}[theo]{Lemma}
\newtheorem{defi}[theo]{Definition}
\newtheorem{prop}[theo]{Proposition}
\newtheorem{cond}[theo]{Condition}
\newtheorem{conj}[theo]{Conjecture}
\newtheorem{cor}[theo]{Corollary}
\newtheorem{remark}[theo]{Remark}
\newtheorem{property}[theo]{Property}
\newtheorem{example}[theo]{Example}
\newtheorem{ques}[theo]{Question}
\numberwithin{equation}{section}

\mathchardef\mhyphen="2D

\def\CP{\mathbb{CP}}
\def\A{{\mathbb A}}
\def\bbI{{\mathbb I}}
\def\N{\mathbb{N}}
\def\bL{\mathbb{L}}
\def\bK{\mathbb{K}}
\def\R{\mathbb{R}}
\def\bS{\mathbb{S}}
\def\C{\mathbb{C}}
\def\Z{\mathbb{Z}}
%\def\bbF{\mathbb{F}}
\def\Q{\mathbb{Q}}
\def\bG{\mathbb{G}}
\def\FS{\mathrm{Lag_{vc}}}
\raggedbottom
\bibcite{...}{...} % so we can \cite{...}

\def\db#1{ \bD^b({#1})}
\def\coh{\operatorname{coh}}
\def\lto{\longrightarrow}
\def\wt{\widetilde}


\def\B{{\mathcal B}}
\def\E{{\mathcal E}}
\def\F{{\mathcal F}}
\def\G{{\mathcal G}}
\def\O{{\mathcal O}}
\def\P{{\mathcal{P}}}
\def\D{{\mathcal{D}}}



\def\bH{{\mathbf H}}
\def\bC{{\mathbf C}}
\def\bD{{\mathbf D}}
\def\bR{{\mathbf R}}
\def\bL{{\mathbf L}}
\def\bY{{\mathbf Y}}

\def\PP{{\mathbb P}}
\def\FF{{\mathbb F}}

\def\pre-tr{\operatorname{pre-tr}}
\def\h{\operatorname{h}}
\def\Maps{\operatorname{Maps}}
\def\Hom{\operatorname{Hom}}
\def\Map{\operatorname{Map}}
\def\End{\operatorname{End}}
\def\NS{\operatorname{NS}}
\def\Pic{\operatorname{Pic}}
\def\rad{\operatorname{rad}}
\def\gr{\operatorname{gr}}
\def\hd{\operatorname{hd}}



\DeclareMathOperator*{\colim}{colim}
\DeclareMathOperator{\colimm}{colim}
\DeclareMathOperator*{\hocolim}{hocolim}
\DeclareMathOperator*{\holim}{holim}



\newcommand{\tens}[1]{%
  \mathbin{\mathop{\otimes}\displaylimits_{#1}}%
}

\newcommand{\Ltens}[1]{%
  \mathbin{\mathop{\otimes}\displaylimits^{\bL}_{#1}}%
}

\newcommand{\conttens}[1]{%
  \mathbin{\mathop{\otimes}\displaylimits^{\cont}_{#1}}%
}

\newcommand{\laxtimes}[1]{\mathop{\times\mkern-13mu\raise1.3ex\hbox{$\scriptscriptstyle\to$}_{#1}}}

\newcommand{\ms}[1]{\mathscr{#1}}
\newcommand{\msK}{\ms{K}}

\newcommand{\hy}{\mhyphen}

\newcommand{\indlim}[1][]{\mathop{\varinjlim}\limits_{#1}}
\newcommand{\inddlim}[1][]{{``{\indlim[#1]}"}}
\newcommand{\prolim}[1][]{\mathop{\varprojlim}\limits_{#1}}
\newcommand{\proolim}[1][]{{``{\prolim[#1]}"}}
\newcommand{\biggplus}[1][]{\mathop{\bigoplus}\limits_{#1}}
\newcommand{\bigooplus}[1][]{{``{\biggplus[#1]}"}}
\newcommand{\prodd}[1][]{\mathop{\prod}\limits_{#1}}
\newcommand{\coprodd}[1][]{\mathop{\coprod}\limits_{#1}}

\newcommand{\limproj}{\varprojlim}
\newcommand{\liminj}{\varinjlim}
\newcommand{\lcone}{\triangleleft}
\newcommand{\rcone}{\triangleright}
\newcommand{\bbar}{\overline}
\newcommand{\hhat}{\widehat}
\newcommand{\toto}{\rightrightarrows}
\newcommand{\xto}{\xrightarrow}
\newcommand{\xlto}{\xleftarrow}
\newcommand{\hto}{\hookrightarrow}
\newcommand{\xhto}{\xhookrightarrow}
\newcommand{\onto}{\twoheadrightarrow}
\newcommand{\xonto}{\xtwoheadrightarrow}
\newcommand{\leftto}{\leftarrow}
\newcommand{\waveto}{\rightsquigarrow}
\newcommand{\modmod}{/ \! \! /}
\newcommand{\bnd}{\operatorname{bnd}}
\newcommand{\lex}{\operatorname{lex}}
\newcommand{\rex}{\operatorname{rex}}
\newcommand{\rig}{\operatorname{rig}}
\newcommand{\oplax}{\operatorname{oplax}}
\newcommand{\lax}{\operatorname{lax}}
\newcommand{\fin}{\operatorname{Fin}}
\newcommand{\profin}{\operatorname{ProFin}}
\newcommand{\comphaus}{\operatorname{CompHaus}}
\newcommand{\Tate}{\operatorname{Tate}}
\newcommand{\Mot}{\operatorname{Mot}}
\newcommand{\Corr}{\operatorname{Corr}}
\newcommand{\QCoh}{\operatorname{QCoh}}
\newcommand{\Nuc}{\operatorname{Nuc}}
\newcommand{\Solid}{\operatorname{Solid}}
\newcommand{\Coh}{\operatorname{Coh}}
\newcommand{\IndCoh}{\operatorname{IndCoh}}
\newcommand{\Proj}{\operatorname{Proj}}
\newcommand{\Flat}{\operatorname{Flat}}
\newcommand{\fla}{\operatorname{flat}}
\newcommand{\Cat}{\operatorname{Cat}}
\newcommand{\Ring}{\operatorname{Ring}}
\newcommand{\Cone}{\operatorname{Cone}}
\newcommand{\Fiber}{\operatorname{Fiber}}
\newcommand{\Cofib}{\operatorname{Cofib}}
\newcommand{\el}{\operatorname{el}}
\newcommand{\idem}{\operatorname{idem}}
\newcommand{\prim}{\operatorname{prim}}
\newcommand{\loc}{\operatorname{loc}}
\newcommand{\cont}{\operatorname{cont}}
\newcommand{\red}{\operatorname{red}}
\newcommand{\acc}{\operatorname{acc}}
\newcommand{\ex}{\operatorname{ex}}
\newcommand{\strcont}{\operatorname{strcont}}
\newcommand{\St}{\operatorname{St}}
\newcommand{\st}{\operatorname{st}}
\newcommand{\cg}{\operatorname{cg}}
\newcommand{\Prr}{\operatorname{Pr}}
\newcommand{\dual}{\operatorname{dual}}
\newcommand{\Hoch}{\operatorname{Hoch}}
\newcommand{\Pol}{\operatorname{Pol}}
\newcommand{\pol}{\operatorname{pol}}
\newcommand{\deff}{\operatorname{def}}
\newcommand{\Fil}{\operatorname{Fil}}
\newcommand{\bbA}{{\mathbb A}}
\newcommand{\bbC}{{\mathbb C}}
\newcommand{\bbH}{{\mathbb H}}
\newcommand{\bbL}{{\mathbb L}}
\newcommand{\bbR}{{\mathbb R}}
\newcommand{\bbW}{{\mathbb W}}
\newcommand{\bbZ}{{\mathbb Z}}
\newcommand{\bbD}{{\mathbb D}}
\newcommand{\bbP}{{\mathbb P}}
\newcommand{\bbN}{{\mathbb N}}
\newcommand{\bE}{{\mathbb E}}
\newcommand{\bbF}{{\mathbb F}}
\newcommand{\bbG}{{\mathbb G}}
\newcommand{\bbQ}{{\mathbb Q}}

\newcommand{\mk}{\mathrm k}
\newcommand{\ev}{\mathrm ev}


\newcommand{\ze}{\zeta _{\sigma}}

\newcommand{\cJ}{{\mathcal J}}
\newcommand{\cQ}{{\mathcal Q}}
\newcommand{\cF}{{\mathcal F}}
\newcommand{\cG}{{\mathcal G}}
\newcommand{\cO}{{\mathcal O}}
\newcommand{\cP}{{\mathcal P}}
\newcommand{\cL}{{\mathcal L}}
\newcommand{\cM}{{\mathcal M}}
\newcommand{\cN}{{\mathcal N}}
\newcommand{\cD}{{\mathcal D}}
\newcommand{\cV}{{\mathcal V}}
\newcommand{\cA}{{\mathcal A}}
\newcommand{\cB}{{\mathcal B}}
\newcommand{\cI}{{\mathcal I}}
\newcommand{\cC}{{\mathcal C}}
\newcommand{\cE}{{\mathcal E}}
\newcommand{\cW}{{\mathcal W}}
\newcommand{\cY}{{\mathcal Y}}
\newcommand{\cU}{{\mathcal U}}
\newcommand{\cR}{{\mathcal R}}
\newcommand{\cS}{{\mathcal S}}
\newcommand{\cT}{{\mathcal T}}
\newcommand{\cK}{{\mathcal K}}
\newcommand{\cH}{{\mathcal H}}
\newcommand{\cX}{{\mathcal X}}
\newcommand{\V}{{\mathcal V}}
\newcommand{\mg}{\mathfrak g}
\newcommand{\mh}{\mathfrak h}
\newcommand{\cHom}{\mathcal Hom}
\newcommand{\mD}{{\mathfrak D}}
\newcommand{\veps}{\varepsilon}

\newcommand{\un}{\underline}
\newcommand{\la}{\langle}
\newcommand{\ra}{\rangle}

\newcommand{\incl}{\operatorname{incl}}
\newcommand{\Compass}{\operatorname{CompAss}}
\newcommand{\prenorm}{\operatorname{PreNorm}}
\newcommand{\Ban}{\operatorname{Ban}}
\newcommand{\Coeq}{\operatorname{Coeq}}
\newcommand{\Eq}{\operatorname{Eq}}
\newcommand{\LEq}{\operatorname{LEq}}
\newcommand{\ann}{\operatorname{ann}}
\newcommand{\cl}{\operatorname{cl}}
\newcommand{\mrB}{\operatorname{B}}
\newcommand{\mrE}{\operatorname{E}}
\newcommand{\qu}{\operatorname{qu}}
\newcommand{\cross}{\operatorname{cr}}
\newcommand{\THH}{\operatorname{THH}}
\newcommand{\HH}{\operatorname{HH}}
\newcommand{\THC}{\operatorname{THC}}
\newcommand{\CAlg}{\operatorname{CAlg}}
\newcommand{\DGAlg}{\operatorname{DGAlg}}
\newcommand{\U}{\operatorname{U}}
\newcommand{\DG}{\operatorname{DG}}
\newcommand{\Fun}{\operatorname{Fun}}
\newcommand{\Tot}{\operatorname{Tot}}
\newcommand{\topp}{\operatorname{top}}
\newcommand{\Lim}{\operatorname{Lim}}	
\newcommand{\Mix}{\operatorname{Mix}}
\newcommand{\Sets}{\operatorname{Sets}}
\newcommand{\Ab}{\operatorname{Ab}}
\newcommand{\Def}{\operatorname{Def}}
\newcommand{\Add}{\operatorname{Add}}
\newcommand{\Forg}{\operatorname{Forg}}
\newcommand{\Perf}{\operatorname{Perf}}
\newcommand{\PsPerf}{\operatorname{PsPerf}}
\newcommand{\perf}{\operatorname{perf}}
\newcommand{\pspe}{\operatorname{pspe}}
\newcommand{\propp}{\operatorname{prop}}
\newcommand{\Perv}{\operatorname{Perv}}
\newcommand{\Kar}{\operatorname{Kar}}
\newcommand{\Kos}{\operatorname{Kos}}
\newcommand{\rat}{\operatorname{rat}}
\newcommand{\PSh}{\operatorname{PSh}}
\newcommand{\copsh}{\operatorname{coPSh}}
\newcommand{\Cosh}{\operatorname{Cosh}}
\newcommand{\Sh}{\operatorname{Shv}}
\newcommand{\Shv}{\operatorname{Shv}}
\newcommand{\supp}{\operatorname{Supp}}
\newcommand{\Ker}{\operatorname{Ker}}
\newcommand{\coker}{\operatorname{coker}}
\newcommand{\Mor}{\operatorname{Mor}}
\newcommand{\im}{\operatorname{Im}}
\newcommand{\re}{\operatorname{Re}}
\newcommand{\codim}{\operatorname{codim}}
\newcommand{\cha}{\operatorname{char}}
\newcommand{\ch}{\operatorname{ch}}
\newcommand{\rank}{\operatorname{rank}}
\newcommand{\rk}{\operatorname{rk}}
\newcommand{\Diff}{\operatorname{Diff}}
\newcommand{\Rep}{\operatorname{Rep}}
\newcommand{\Crit}{\operatorname{Crit}}
\newcommand{\Nilp}{\operatorname{Nilp}}
\newcommand{\Mat}{\operatorname{Mat}}
\newcommand{\Eu}{\operatorname{Eu}}
\newcommand{\eu}{\operatorname{eu}}
\newcommand{\Ext}{\operatorname{Ext}}
\newcommand{\Tor}{\operatorname{Tor}}
\newcommand{\isomo}{\overset{\sim}{=}}
\newcommand{\isomoto}{\overset{\sim}{\to}}
\newcommand{\Id}{\operatorname{Id}}
\newcommand{\tr}{\operatorname{tr}}
\newcommand{\Tr}{\operatorname{Tr}}
\newcommand{\str}{\operatorname{str}}
\newcommand{\Nm}{\operatorname{Nm}}
\newcommand{\shHom}{\underline{\operatorname{Hom}}}
\newcommand{\Symm}{\operatorname{Sym}}
\newcommand{\Arg}{\operatorname{Arg}}
\newcommand{\bbPr}{\operatorname{pr}}
\newcommand{\Star}{\operatorname{Star}}
\newcommand{\Sp}{\operatorname{Sp}}
\newcommand{\Span}{\operatorname{Span}}
\newcommand{\add}{\operatorname{add}}
\newcommand{\smd}{\operatorname{smd}}
\newcommand{\res}{\operatorname{res}}
\newcommand{\cores}{\operatorname{cores}}
\newcommand{\Or}{{\underline{o}}}
\newcommand{\Aut}{\operatorname{Aut}}
\newcommand{\rmMod}{\operatorname{Mod}}
\newcommand{\ord}{\operatorname{ord}}
\newcommand{\even}{\operatorname{even}}
\newcommand{\odd}{\operatorname{odd}}
\newcommand{\Sym}{\operatorname{Sym}}
\newcommand{\Alt}{\operatorname{Alt}}
\newcommand{\Ind}{\operatorname{Ind}}
\newcommand{\Pro}{\operatorname{Pro}}
\newcommand{\Calk}{\operatorname{Calk}}
\newcommand{\Res}{\operatorname{Res}}
\newcommand{\Spec}{\operatorname{Spec}}
\newcommand{\Spf}{\operatorname{Spf}}
\newcommand{\Hilb}{\rm Hilb}
\newcommand{\Stab}{\rm Stab}
\newcommand{\Sign}{\rm Sign}
\newcommand{\sgn}{\rm sgn}
\newcommand{\Vect}{{\rm Vect}}
\newcommand{\groth}[1]{K_0[\V_{#1}]}
\newcommand{\grsig}[1]{K_0[\V_{#1}]_\sigma}
\newcommand{\dgart}{\operatorname{dgart}}
\newcommand{\art}{\operatorname{art}}
\newcommand{\Lie}{\operatorname{Lie}}
\newcommand{\coDef}{\operatorname{coDef}}
\newcommand{\cart}{\operatorname{cart}}
\newcommand{\Alg}{\operatorname{Alg}}
\newcommand{\ADGalg}{\operatorname{ADGalg}}
\newcommand{\Grp}{\operatorname{Grp}}
\newcommand{\Ho}{\operatorname{Ho}}
\newcommand{\SF}{\operatorname{SF}}
\newcommand{\id}{\operatorname{id}}
\newcommand{\QP}{\operatorname{QP}}
\newcommand{\dgalg}{\operatorname{dgalg}}
\newcommand{\dgcat}{\operatorname{dgcat}}
\newcommand{\DEF}{\operatorname{DEF}}
\newcommand{\coDEF}{\operatorname{coDEF}}
\newcommand{\eval}{\operatorname{ev}}
\newcommand{\coev}{\operatorname{coev}}
\newcommand{\Tors}{\operatorname{Tors}}
\newcommand{\tors}{\operatorname{tors}}
\newcommand{\QGr}{\operatorname{QGr}}
\newcommand{\Comp}{\operatorname{Comp}}
\newcommand{\comp}{\operatorname{comp}}
\newcommand{\compl}{\operatorname{compl}}
\newcommand{\QMod}{\operatorname{QMod}}
\newcommand{\Acycl}{\operatorname{Acycl}}
\newcommand{\acycl}{\operatorname{acycl}}
\newcommand{\qgr}{\operatorname{qgr}}
\newcommand{\qmod}{\operatorname{qmod}}
\newcommand{\adgalg}{\operatorname{adgalg}}
\newcommand{\Set}{\operatorname{Set}}
\newcommand{\sSet}{\operatorname{sSet}}
\newcommand{\alg}{\operatorname{alg}}
\newcommand{\calg}{\operatorname{calg}}
\newcommand{\ad}{\operatorname{ad}}
\newcommand{\NGr}{\operatorname{NGr}}
\newcommand{\Gr}{\operatorname{Gr}}
\newcommand{\Fl}{\operatorname{Fl}}
\newcommand{\pr}{\operatorname{pr}}
\newcommand{\pt}{\operatorname{pt}}
\newcommand{\vdim}{\operatorname{vdim}}
\newcommand{\Mod}{\operatorname{Mod}}
\newcommand{\clopen}{\operatorname{ClOpen}}
\newcommand{\Open}{\operatorname{Open}}
\newcommand{\mult}{\operatorname{mult}}
\newcommand{\trdeg}{\operatorname{trdeg}}
\newcommand{\solid}{{\scalebox{0.5}{$\square$}}}
\newcommand{\bs}{\blacksquare}

\usepackage{epsf}
\usepackage{amscd}
\newcommand{\mP}{\mathcal{P}}
\newcommand{\I}{\mathcal{I}}
\newcommand{\T}{\mathcal{T}}
\newcommand{\mR}{\mathcal{R}}
\newcommand{\mQ}{\mathcal{Q}}
\newcommand{\mfp}{\mathfrak{p}}
\newcommand{\mfq}{\mathfrak{q}}
\newcommand{\mS}{\mathfrak{S}}
\newcommand{\M}{\mathcal{M}}
%\newcommand{\G}{\mathcal{G}}
\newcommand{\m}{\mathfrak{m}}
\newcommand{\ma}{\mathfrak{a}}
\newcommand{\mIm}{\mathrm{Im}}
\newcommand{\bfL}{\mathbf{L}}
\newcommand{\bfR}{\mathbf{R}}
\newcommand{\bw}{\mathrm{w}}
\newcommand{\rmmod}{\operatorname{mod}}
\newcommand{\rmcomod}{\operatorname{comod}}
\newcommand{\one}{\mathbf{1}}
\newcommand{\rarr}{\rightarrow}
\renewcommand{\theequation}{\thesection.\arabic{equation}}

\title[Some remarks on Quillen's D\'evissage theorem]
{Some remarks on Quillen's D\'evissage theorem}

\author{Alexander I. Efimov}
\address{Steklov Mathematical Institute of RAS, Gubkin St. 8, GSP-1, Moscow 119991, Russia}
\email{efimov@mccme.ru}

%\thanks{The author was partially supported by the European Research Council (ERC, CurveArithmetic, 101078157).}


%\thanks{The author is partially supported by Laboratory of Mirror Symmetry NRU HSE, RF government  grant, ag. N 14.641.31.0001}
%\thanks{MSC: 16E40, 18G40, 58K05, 11R04.}
%\thanks{The author was partially supported by RFBR, grant 2998.2014.1, and RFBR, research project 15-51-50045.}

\begin{document}

\begin{abstract} In this paper we give a different proof of Quillen's D\'evissage theorem using Barwick's theorem of the heart. The key ingredient is a certain short exact sequence of dg categories, which is closely related with the Auslander-type construction for nilpotent extensions which was used in the papers of Kuznetsov-Lunts \cite{KL15}, Land-Tamme \cite{LT19} and the author \cite{E20}.
\end{abstract}

%\keywords{Mac Lane homology and cohomology; Wieferich primes; critical points; spectral sequences; cubical construction.}

\maketitle

\tableofcontents

\section{Introduction}
\label{sec:intro}

We recall the following classical result of Quillen.

\begin{theo}\label{th:Quillen} \cite[Theorem 4]{Qui}
Let $\cA$ be a small abelian category, and let $\cB\subset\cA$ be a strictly full abelian subcategory which is closed under taking subobjects and quotients. Suppose that each object of $\cA$ has a finite filtration with subquotients in $\cB.$ Then we have isomorphisms $K_n(\cB)\xto{\sim}K_n(\cA)$ for $n\geq 0.$
\end{theo}

Another important result which is due to Barwick is about $K$-theory of stable $\infty$-categories with bounded $t$-structures.

\begin{theo}\label{th:Barwick} \cite[Theorem 6.1]{Bar15}
Let $\cT$ be a small stable $\infty$-category with a bounded $t$-structure, i.e. $\cT$ is generated as a stable subcategory by the heart $\cT^{\heartsuit}.$ Then we have isomorphisms $K_n(\cT^{\heartsuit})\xto{\sim} K_n(\cT)$ for $n\geq 0.$
\end{theo}
 
Both results are extremely useful for computations of $K$-theory. The original proof of Theorem \ref{th:Quillen} in \cite{Qui} uses Quillen's Theorem A and the $Q$-construction, and the argument is quite mysterious from the categorical perspective. On the other hand, Barwick's proof of Theorem \ref{th:Barwick} uses the machinery of exact $\infty$-categories as defined in \cite{Bar15}, and the argument is much more direct. In this paper we show that in fact Theorem \ref{th:Barwick} implies Theorem \ref{th:Quillen}.

We now explain the main construction. Let $\cA\supset \cB$ be as in Theorem \ref{th:Quillen}. Suppose that moreover each object of $\cA$ is an extension of two objects of $\cB.$ Note that the general case reduces to this case.

Consider the category $\cE_{\cA,\cB}$ of pairs $(X,Y),$ where $X\in\cA$ and $Y\subset X$ is a subobject such that $Y\in\cB$ and $X/Y\in\cB.$ Then the category $\cE_{\cA,\cB}$ is quasi-abelian, in particular we can consider it as an exact category. We refer to Subsection \ref{ssec:quasi_abelian} for the definition of a quasi-abelian category. 

We prove the following result.

\begin{theo}\label{th:key_construction_intro}
Within the above notation, the following holds.
\begin{enumerate}[label=(\roman*), ref=(\roman*)]
	\item We have a natural semi-orthogonal decomposition
	\begin{equation*}
		D^b(\cE_{\cA,\cB})=\la D^b(\cB), D^b(\cB)\ra.
	\end{equation*}
	We denote by $\Phi_1,\Phi_2:D^b(\cB)\to D^b(\cE_{\cA,\cB})$ the inclusion functors for the $1$-st resp. $2$-nd component. We denote by $\Phi_1^L$ resp. $\Phi_2^R$ the left resp. right adjoint to $\Phi_1$ resp. $\Phi_2.$ \label{SOD}
	\item We have a short exact sequence of dg categories
	\begin{equation}\label{eq:key_ses}
		0\to \cT_{\cA,\cB}\xto{\Psi} D^b(\cE_{\cA,\cB})\xto{q} D^b(\cA)\to 0.
	\end{equation}
	Here the functor $q$ is induced by the (exact) forgetful functor $\cE_{\cA,\cB}\to \cA.$ The functor $\B\to \cT_{\cA,\cB},$ $X\mapsto\Cone((X,0)\to (X,X)),$ is fully faithful, and its essential image is the heart of a bounded $t$-structure on $\cT_{\cA,\cB}.$ \label{ses}
	\item Both compositions $q\circ \Phi_i,$ $i=1,2,$ are isomorphic to the derived functor of the inclusion $\cB\hto\cA.$ The compositions $\Phi_1^L[-1]\circ \Psi$ and $\Phi_2^R\circ \Psi$ are left inverses to the realization functor $D^b(\cB)\to \cT_{\cA,\cB}.$ \label{functors_on_components}
\end{enumerate}
\end{theo}

Theorem \ref{th:Quillen} follows almost immediately from Theorems \ref{th:key_construction_intro} and \ref{th:Barwick}. 
  
Part \ref{SOD} of Theorem \ref{th:key_construction_intro} is straightforward. To prove \ref{ses} and \ref{functors_on_components} we give an explicit description of the category $\cL\cH(\cE_{\cA,\cB})$ -- the right abelian envelope of the category $\cE_{\cA,\cB},$ see Section \ref{sec:description_of_the_envelope}. Recall from \cite{Sch99} that for a quasi-abelian category $\cE$ its right abelian envelope is the heart of the bounded $t$-structure on $D^b(\cE)$ whose connective part is the subcategory of complexes in non-negative homological degrees. This heart consists of complexes of the form $\Cone(X\xto{f} Y),$ where $f$ is a monomorphism (not necessarily strict), see Proposition \ref{prop:generalities_on_abelian_envelope}. Both functors in \eqref{eq:key_ses} are $t$-exact, where we take the standard $t$-structure on $D^b(\cA).$

The category $D^b(\cE_{\cA,\cB})$ was in fact considered previously (with a different description) in \cite{KL15, E20} in the case $\cA=\Coh(A),$ $\cB=\Coh(A/I),$ where $A$ is a commutative noetherian ring and $I\subset A$ is an ideal such that $I^2=0.$ This category is also closely related with a certain lax pullback which was implicitly used in \cite[Proofs of Theorem 2.25 and Corollary 3.5]{LT19}, see Example \ref{ex:familiar_example} for details.

{\noindent{\bf Convention.}} We will use only the homological grading.

{\noindent{\bf Acknowledgements.}} This work was performed at the Steklov International Mathematical Center
and supported by the Ministry of Science and Higher Education of the
Russian Federation (agreement no. 075-15-2025-303). 

\section{Preliminaries}

\subsection{Quasi-abelian categories}
\label{ssec:quasi_abelian}

We refer to \cite{Sch99} for the general theory of quasi-abelian categories. The closely related notion is that of a torsion theory in an abelian category as defined in \cite{HRS96}. This relation is studied in \cite[Section 5.4 and Appendix B]{BVdB03}. 

Recall that an additive category $\cE$ is called quasi-abelian if $\cE$ has kernels and cokernels, the class of strict monomorphisms is closed under pushouts and the class of strict epimorphisms is closed under pullbacks. Here a monomorphism $f$ is strict if $f=\ker(\coker(f)),$ and similarly for epimorphisms. In particular, a quasi-abelian category $\cE$ has an exact structure, for which the admissible short exact sequences are of the form
\begin{equation}\label{eq:adm_ses}
0\to X\xto{f} Y\xto{g} Z\to 0,
\end{equation}
where $f=\ker(g)$ and $g=\coker(f).$ In particular, we have the bounded derived category $D^b(\cE),$ which we consider as a dg category over $\Z.$

Clearly, any abelian category is quasi-abelian. An additive functor between quasi-abelian categories is called exact if it preserves admissible short exact sequences.
A functor $F:\cE\to\cA$ from a quasi-abelian category $\cE$ to an abelian category $\cA$ is called right exact if for any admissible short exact sequence \eqref{eq:adm_ses} we have an exact sequence
\begin{equation*}
F(X)\to F(Y)\to F(Z)\to 0.
\end{equation*}

\begin{prop}\label{prop:generalities_on_abelian_envelope}\cite[Definition 1.2.18, Propositions 1.2.32 and 1.2.34]{Sch99} Let $\cE$ be a small quasi-abelian category. 
\begin{enumerate}[label=(\roman*),ref=(\roman*)]
	\item The category $D^b(\cE)$ has a bounded $t$-structure, whose heart consists of objects of the form $\Cone(X\xto{f} Y),$ where $f$ is a monomorphism (but not necessarily a strict monomorphism). This heart is denoted by $\cL\cH(\cE).$ The inclusion functor $\cE\to\cL\cH(\cE),$ $X\mapsto X,$ is exact.
	\item The functor $D^b(\cE)\to D^b(\cL\cH(\cE))$ is an equivalence.
	\item The functor $\cE\to \cL\cH(\cE)$ satisfies the following universal property: for any small abelian category $\cA$ it induces an equivalence
	\begin{equation*}
	\Fun^{\rex}(\cL\cH(\cE),\cA)\xto{\sim}\Fun^{\rex}(\cE,\cA).
	\end{equation*}
	Here the superscript $\rex$ means ``right exact''.
\end{enumerate}
\end{prop}

Following \cite[Section 1.2.4]{Sch99}, we call the category $\cL\cH(\cE)$ the (right) abelian envelope of $\cE.$ We recall the following characterization of the category $\cL\cH(\cE).$

\begin{prop}\label{prop:characterization_of_LH}\cite[Proposition 1.2.36]{Sch99}
	Let  $\cE$ be a small quasi-abelian category and let $\cA$ be a small abelian category. Suppose that we have a fully faithful functor $F:\cE\to\cA$ such that the following conditions hold.
	\begin{enumerate}[label=(\roman*),ref=(\roman*)]
		\item The essential image $F(\cE)\subset\cA$ is closed under taking subobjects. 
		\item For every object $X\in\cA$ there exists an object $Y\in\cE$ and an epimorphism $F(Y)\onto X.$
	\end{enumerate}
	Then $F$ extends to an equivalence of categories $\cL\cH(\cE)\xto{\sim}\cA.$
\end{prop}

For completeness we also recall the relation between quasi-abelian categories and torsion pairs in abelian categories. Recall from \cite{HRS96} that for an abelian category $\cA$ a pair of full subcategories $(\cT,\cF)$ is called a torsion pair if $\Hom(\cT,\cF)=0$ and for any $X\in\cA$ there is a short exact sequence
\begin{equation*}
0\to Y\to X\to Z\to 0,\quad Y\in\cT,\,Z\in\cF.
\end{equation*}
In this case $\cT$ and $\cF$ are additive subcategories closed under extensions, $\cT$ is closed under taking quotients and $\cF$ is closed under taking subobjects. Moreover, both $\cT$ and $\cF$ are quasi-abelian and the inclusion functors $\cT\to\cA$ and $\cF\to\cA$ are exact. Such a torsion pair is called cotilting if every object of $\cA$ is a quotient of an object of $\cF.$

\begin{prop}\cite[Proposition B.3]{BVdB03}
Let $\cE$ be a small additive category. The following are equivalent.
\begin{enumerate}[label=(\roman*),ref=(\roman*)]
	\item $\cE$ is quasi-abelian.
	\item There exists a cotilting torsion pair $(\cT,\cF)$ in a small abelian category $\cA$ with an equivalence $\cE\simeq \cF.$ \label{via_cotilting_pair}
\end{enumerate}
In the situation of \ref{via_cotilting_pair} we have $\cA\simeq\cL\cH(\cE).$
\end{prop} 

\subsection{Ind-objects and pro-objects}
\label{ssec:ind_obj_pro_obj}

We refer to \cite{KS01} for the general theory of ind-objects in ordinary categories. For a small abelian category $\cA$ we have the abelian category $\Ind(\cA)$ of ind-objects, and the inclusion functor $\cA\to\Ind(\cA)$ is exact. We write the objects of $\Ind(\cA)$ in the form $\inddlim[i\in I]X_i,$ where $I$ is a directed poset and we are assuming a functor $I\to\cA,$ $i\mapsto X_i.$ Recall that we have an equivalence
\begin{equation*}
\Ind(\cA)\simeq \Fun^{\lex}(\cA^{op},\Ab),\quad \inddlim[i] X_i\mapsto \indlim[i] \h_{X_i},
\end{equation*}
where $h_{X_i}=\Hom(-,X_i).$ Here the superscript $\lex$ means ``left exact''. 

It is convenient to give another equivalent description of the category $\Ind(\cA)$ via sheaves. Namely, consider the Grothendieck topology on the category $\cA,$ for which a sieve $S\subset h_X$ is a covering sieve if it contains an epimorphism $Y\onto X.$ We will call it the single epi topology, following \cite{Ka22}. Then a functor $F:\cA^{op}\to \Ab$ is left exact if and only if $F$ is a sheaf for the single epi topology, which is also additive as a functor. Hence, we have an equivalence
\begin{equation*}
\Ind(\cA)\simeq \Shv^{\add}(\cA;\Ab).
\end{equation*}

Given a functor $F:\cA^{op}\to\Ab,$ considered as a presheaf of abelian groups, we denote by $F^{\sharp}$ its sheafification for the single epi topology.  We will use the following standard statement.

\begin{prop}\label{prop:sheafification_of_additive_presheaves}
For an additive functor $F:\cA^{op}\to \Ab,$ the sheafification $F^{\sharp}$ is an additive sheaf. In particular, the inclusion $\Ind(\cA)\to \Fun^{\add}(\cA^{op},\Ab)$ has an exact left adjoint, given by the composition
\begin{equation}\label{eq:left_adjoint_via_sheafification}
\Fun^{\add}(\cA^{op},\Ab)\xto{(-)^{\sharp}} \Shv^{\add}(\cA^{op},\Ab)\simeq \Ind(\cA).
\end{equation} 
\end{prop}

\begin{proof}
We only need the following observation: for any objects $X,Y\in\cA$ and an epimorphism $(f,g):Z\onto X\oplus Y,$ we can find a pair of epimorphisms $f':X'\onto X$ and $g':Y'\onto Y$ such that the morphism $(f',g'):X'\oplus Y'\to X\oplus Y$ factors through $Z.$ Namely, it suffices to take $X'=\ker(g),$ $Y'=\ker(f),$ and define $f'$ and $g'$ as the compositions $X'\to Z\xto{f}X$ and $Y'\to Z\xto{g}Y.$

Exactness of the functor \eqref{eq:left_adjoint_via_sheafification} is automatic since the sheafification is always exact.
\end{proof}

Note that exactness of the functor \eqref{eq:left_adjoint_via_sheafification} is also a trivial special case of the Gabriel-Kuhn-Popesco theorem \cite{GP64, Ku94}.

We also have the abelian category of pro-objects $\Pro(\cA)\simeq \Ind(\cA^{op})^{op}.$ We write the objects of $\Pro(\cA)$ in the form $\proolim[i\in I]X_i,$ where $I$ is a codirected poset and we are assuming a functor $I\to\cA,$ $i\mapsto X_i.$ We will need the following immediate application of Proposition \ref{prop:sheafification_of_additive_presheaves}.

\begin{cor}\label{cor:making_a_functor_right_exact}
Let $\cA$ and $\cB$ be small abelian categories. Then the inclusion functor
\begin{equation*}
\incl:\Fun^{\rex}(\cA,\Pro(\cB))\to \Fun^{\add}(\cA,\Pro(\cB))
\end{equation*}
has a right adjoint $\incl^R.$ Given an additive functor $F:\cA\to\Pro(\cB)$ and an object $X\in\cB,$ the morphism of functors
\begin{equation*}
\Hom(F(-),X)\to \Hom(\incl^R(F)(-),X)
\end{equation*}
exhibits the target as a sheafification of the source for the single epi topology on $\cA.$ 
\end{cor} 

\begin{proof}
Indeed, we have a commutative square
\begin{equation*}
\begin{CD}
\Fun^{\rex}(\cA,\Pro(\cB)) @>{\sim}>> \Fun^{\lex}(\cB,\Shv^{\add}(\cA;\Ab))^{op}\\
@V{\incl}VV @VVV\\
\Fun^{\add}(\cA,\Pro(\cB)) @>{\sim}>> \Fun^{\lex}(\cB,\Fun^{\add}(\cA^{op},\Ab))^{op}.
\end{CD}
\end{equation*}
Here the horizontal functors are given by $F\mapsto (X\mapsto \Hom(F(-),X)).$ The corollary now follows directly from Proposition \ref{prop:sheafification_of_additive_presheaves}.
\end{proof}

\section{Right abelian envelope of the category $\cE_{\cA,\cB}$}
\label{sec:description_of_the_envelope}

Throughout this section, $\cA\supset\cB$ are small abelian categories as in Theorem \ref{th:key_construction_intro}. That is, $\cB$ is a strictly full subcategory of $\cA,$ which is closed under taking subobjects and quotients, and each object of $\cA$ is an extension of two objects of $\cB.$ The category $\cE_{\cA,\cB}$ is defined as in the introduction, i.e. it is the category of pairs $(X,Y),$ where $X\in\cA,$ $Y\subset X$ and $Y\in\cB,$ $X/Y\in\cB.$ We will describe directly the right abelian envelope $\cL\cH(\cE_{\cA,\cB})$ and deduce some of its structural properties.

We start with the following trivial observation.

\begin{lemma}\label{lem:intersections}
Let $X\in \cA$ be an object and suppose that $Y_1,Y_2\subset X$ are subobjects such that $X/Y_1, X/Y_2\in \cB.$ Then $X/(Y_1\cap Y_2)\in\cB.$
\end{lemma}

\begin{proof}
Indeed, this follows from the monomorphism $X/(Y_1\cap Y_2)\to X/Y_1\oplus X/Y_2.$
\end{proof}

We denote by $i:\cB\to\cA$ the inclusion functor, and abusing the notation we denote by the same symbol the induced functor on the pro-completions $i:\Pro(\cB)\to\Pro(\cA)$.
The latter functor has a left adjoint $i^L:\Pro(\cA)\to\Pro(\cB),$ which (commutes with cofiltered limits and) on the objects of $\cA$ is given by
\begin{equation*}
i^L(X) = \proolim[\substack{Y\subset X,\\ Y\in\cB, X/Y\in\cB}] X/Y,\quad X\in\cA.
\end{equation*}
Consider now the functor
\begin{equation}\label{eq:direct_definition_of_j}
j:\cA\to\Pro(\cB),\quad j(X) = \proolim[\substack{Y\subset X,\\ Y\in\cB, X/Y\in\cB}] Y.
\end{equation}
Equivalently, $j$ can be defined by the functorial short exact sequence
\begin{equation}\label{eq:ses_defining_j}
0\to i(j(X))\to X\to i(i^L(X))\to 0.
\end{equation}
The following proposition is immediate from \eqref{eq:ses_defining_j} and also follows directly from \eqref{eq:direct_definition_of_j}. 

\begin{prop}\label{prop:j_preserves_epi}
The functor $j:\cA\to\Pro(\cB)$ preserves epimorphisms.
\end{prop} 

\begin{proof}
Consider a short exact sequence
\begin{equation*}
0\to X\to Y\to Z\to 0
\end{equation*}
in $\cA.$ Applying \eqref{eq:ses_defining_j} to $Y$ and $Z,$ we obtain an exact sequence
\begin{equation*}
X\to \ker(i(i^L(Y))\to i(i^L(Z)))\to \coker(i(j(Y))\to i(j(Z)))\to 0.
\end{equation*}
The first morphism is a composition of epimorphisms
\begin{equation*}
X\to i(i^L(X))\to \ker(i(i^L(Y))\to i(i^L(Z))),
\end{equation*} 
hence it is also an epimorphism. This proves that $j(Y)\to j(Z)$ is an epimorphism.
\end{proof}

Consider the inclusion functor $\incl:\Fun^{\rex}(\cA,\Pro(\cB))\to \Fun^{\add}(\cA,\Pro(\cB)),$ and denote by $\incl^R$ its right adjoint, which is described in Corollary \ref{cor:making_a_functor_right_exact}. We define the right exact functor $\wt{j}:\cA\to \Pro(\cB)$ by $\wt{j}=\incl^R(j).$

\begin{remark}
Another description of the functor $\wt{j}$ is the following. Consider the left derived functor $\bL i^L:D^-(\Pro(\cA))\to D^-(\Pro(\cB)),$ which exists since $\Pro(\cA)$ has enough projective objects. Then for $X\in\cA$ we have an isomorphism
\begin{equation*}
i(\wt{j}(X))\cong H_0(\Fiber(X\to i(\bL i^L(X)))).
\end{equation*}
However, the former description of $\wt{j}$ will be more useful for us.
\end{remark}

We will need the following observation.

\begin{prop}\label{prop:j_tilde_to_j_is_epi}
For $X\in\cA,$ the natural morphism $\wt{j}(X)\to j(X)$ is an epimorphism.
\end{prop}

\begin{proof}
Take some object $Y\in \cB.$ By Proposition \ref{prop:j_preserves_epi}, the functor $\Hom(j(-),Y)$ is a separated presheaf on $\cA$ with respect to the single epi topology. Its sheafification is given by $\Hom(\wt{j}(-),Y)$ by Corollary \ref{cor:making_a_functor_right_exact}. Hence, we have a monomorphism of presheaves $\Hom(j(-),Y)\to \Hom(\wt{j}(-),Y).$ This exactly means that the morphism $\wt{j}(X)\to j(X)$ is an epimorphism for each $X\in\cA.$
\end{proof}

Now we define an abelian category, which will be in fact equivalent to the right abelian envelope $\cL\cH(\cE_{\cA,\cB}).$

{\noindent{\bf Construction.}} {\it We define the category $\cC_{\cA,\cB}$ as follows. Its objects are quadruples $(X,Y,u,v),$ where $X\in\cA,$ $Y\in\cB,$ $u:i(Y)\to X,$ $v:\wt{j}(X)\to Y,$ such that the following conditions hold:
\begin{enumerate}[label=(\roman*),ref=(\roman*)]
	\item The composition $\wt{j}(i(Y))\xto{\wt{j}(u)}\wt{j}(X)\xto{v}Y$ is zero.
	\item The following square commutes:
	\begin{equation}\label{eq:comm_square_for_u_and_v}
		\begin{CD}
			i(\wt{j}(X)) @>{i(v)}>> i(Y)\\
			@VVV @V{u}VV\\
			i(j(X)) @>>> X.
		\end{CD}
	\end{equation}
	Here the left vertical and the lower horizontal maps are the canonical ones. 
\end{enumerate}
The morphisms in $\cC_{\cA,\cB}$ are defined in the natural way.}

To give a feeling of the category $\cC_{\cA,\cB},$ we mention the following special case.

\begin{example}\label{ex:familiar_example}
Let $A$ be a right noetherian associative unital ring with a two-sided ideal $I\subset A$ such that $I^2=0.$ Consider the abelian category $\cA=\Coh\hy A$ of finitely generated right $A$-modules, and the full subcategory $\cB=\Coh\hy A/I\subset \cA.$ Then $\cA$ and $\cB$ satisfy the above conditions. The functors $j,\wt{j}:\cA\to\Pro(\cB)$ take values in the constant pro-objects and are given by
\begin{equation*}
j(M)=MI,\quad \wt{j}(M)=M\tens{A}I.
\end{equation*}
We have $\cC_{\cA,\cB}\simeq \Coh\hy D,$ where the (right noetherian) ring $D$ can be depicted as follows:
\begin{equation*}
D=\begin{pmatrix}
	A & I\\
	A/I & A/I
\end{pmatrix}.
\end{equation*}
This means that $D$ is the direct sum of matrix entries, and $D$ is equipped with the matrix multiplication. The categories $D^b(\Coh\hy D)$ and $\Perf(D)$ were used respectively in \cite[Section 8.1]{E20} and \cite[Section 5]{KL15} in the case when $A$ is commutative (this is a special case of a construction for a general nilpotent ideal).

The category $\Perf(D)$ was also used implicitly in \cite[Proofs of Theorem 2.25 and Corollary 3.5]{LT19}, when $A$ is not necessarily right noetherian. Namely, within the notation of loc.cit., this category is equivalent to the lax pullback:
\begin{equation}\label{eq:equiv_with_lax_pullback}
\Perf(D)\xto{\sim}\Perf(C(I,A))\laxtimes{\Perf(C(I,A/I))} \Perf(A/I).
\end{equation}
Here $C(I,A)$ and $C(I,A/I)$ are dg rings given by
\begin{equation*}
C(I,A)=\Cone(I\hto A),\quad C(I,A/I)=\Cone(I\xto{0}A/I),
\end{equation*}
and the map $C(I,A)\to C(I,A/I)$ is given by the identity on $I$ and by the projection on $A.$ The equivalence \eqref{eq:equiv_with_lax_pullback} is given on the generator by 
\begin{equation*}
D\mapsto (C(I,A),A/I,\id)\oplus (C(I,A),0,0).
\end{equation*}
\end{example}

We first summarize the basic properties of the category $\cC_{\cA,\cB}.$

\begin{prop}\label{prop:properties_of_C_AB}
\begin{enumerate}[label=(\roman*),ref=(\roman*)]
	\item The category $\cC_{\cA,\cB}$ is abelian. The forgetful functors 
	\begin{equation*}
	\cC_{\cA,\cB}\to\cA,\,(X,Y,u,v)\mapsto X,\quad\text{and}\quad \cC_{\cA,\cB}\to\cB,\,(X,Y,u,v)\mapsto Y,
	\end{equation*} are exact. \label{abelian_category}
	\item For any object $(X,Y,u,v)\in\cC_{\cA,\cB},$ the object $X/u(i(Y))$ is contained in $\cB.$ \label{cokernel_in_B}
	\item Consider an object $(X,Y)\in\cE_{\cA,\cB},$ and let $u:i(Y)\to X$ be the inclusion morphism. Then there exists a unique morphism $v:\wt{j}(X)\to Y$ such that the quadruple $(X,Y,u,v)$ is a well-defined object of $\cC_{\cA,\cB}.$ More precisely, $v$ is given by the composition of natural maps:
	\begin{equation}\label{eq:composition_defining_v}
	\wt{j}(X)\to j(X)=\proolim[\substack{Z\subset X,\\ Z\in\cB, X/Z\in\cB}]\to Y
	\end{equation}
	 This gives a fully faithful functor $\alpha:\cE_{\cA,\cB}\to \cC_{\cA,\cB}.$ Its essential image consists of objects $(X',Y',u',v')$ such that $u'$ is a monomorphism. \label{inclusion_of_E_AB}
	 \item We have a well-defined fully faithful functor $\beta:\cB\to\cC_{\cA,\cB},$ given by $\beta(X)=(0,X,0,0).$ We have a torsion pair in $\cC_{\cA,\cB},$ given by $(\beta(\cB),\alpha(\cE_{\cA,\cB})).$ \label{torsion_pair}
\end{enumerate}
\end{prop}

\begin{proof} \ref{abelian_category} follows directly from the right exactness of the functors $i:\cB\to\cA$ and $\wt{j}:\cA\to\Pro(\cB).$

To prove \ref{cokernel_in_B}, denote by $Y'\subset X$ the image of $u:i(Y)\to X,$ in particular $Y'\in B.$ In the commutative square \eqref{eq:comm_square_for_u_and_v} the left vertical arrow is an epimorphism by Proposition \ref{prop:j_tilde_to_j_is_epi}, hence the canonical morphism $i(j(X))\to X$ factors through $Y'.$ This exactly means that $X/Y'\in\cB,$ as required. 

We prove \ref{inclusion_of_E_AB}. Suppose that $v:\wt{j}(X)\to Y$ is some morphism satisfying the required conditions. As above, in the commutative square \eqref{eq:comm_square_for_u_and_v} the left vertical arrow is an epimorphism. Also, the right vertical arrow is a monomorphism by assumption. Hence, there exists a unique diagonal arrow $i(j(X))\to i(Y),$ making both triangle commutative. This shows that $v$ coincides with the composition \eqref{eq:composition_defining_v}. On the other hand, it is clear that this composition satisfies the required conditions (note that $j(Y)=0$). Fully faithfulness of the resulting functor $\alpha:\cE_{\cA,\cB}\to \cC_{\cA,\cB}$ follows directly. The description of the essential image of $\alpha$ follows using \ref{cokernel_in_B}.

It remains to prove \ref{torsion_pair}. The fully faithfulness of $\beta$ is clear, and so is the vanishing of $\Hom(\beta(Z),\alpha(X,Y))$ for $Z\in\cB,$ $(X,Y)\in\cE_{\cA,\cB}.$ The right adjoint $\beta^R$ to $\beta$ is given by $i(\beta^R(X,Y,u,v))=\ker(u).$ The adjunction counit  $\beta(\beta^R(X,Y,u,v))\to (X,Y,u,v)$ is a monomorphism, and its cokernel is contained in the essential image of $\alpha$ by \ref{inclusion_of_E_AB}. This proves \ref{torsion_pair}.
\end{proof}

We now describe the right abelian envelope of $\cE_{\cA,\cB}.$

\begin{prop}\label{prop:abelian_envelope}
The fully faithful exact inclusion $\alpha:\cE_{\cA,\cB}\to \cC_{\cA,\cB}$ induces an equivalence $\cL\cH(\cE_{\cA,\cB})\simeq \cC_{\cA,\cB}.$ 
\end{prop}

\begin{proof}
By Proposition \ref{prop:properties_of_C_AB} \ref{torsion_pair} the category $\alpha(\cE_{\cA,\cB})$ is the full subcategory of torsion-free objects for a torsion pair in $\cC_{\cA,\cB}.$ Hence, by Proposition \ref{prop:characterization_of_LH} we only need to prove that any object of $\cC_{\cA,\cB}$ is a quotient of an object of $\alpha(\cE_{\cA,\cB}).$

Take some object $(X,Y,u,v)\in\cC_{\cA,\cB}.$ As above, put $Y'=\im(u:i(Y)\to X),$ so $Y'$ is an object of $\cB$ and the morphism $i(j(X))\to X$ factors through $Y'.$  Recall that the functor $\Hom(\wt{j}(-),Y)$ resp. $\Hom(\wt{j}(-),Y')$ is the sheafification of the separated presheaf $\Hom(j(-),Y)$ resp. $\Hom(j(-),Y')$ with respect to the single epi topology on $\cA.$ It follows that there exists an epimorphism $f:Z\to X$ and a morphism $w:j(Z)\to Y$ such that the following diagram commutes:
\begin{equation}\label{eq:two_comm_squares}
\begin{tikzcd}
\wt{j}(Z) \ar[d, "\wt{j}(f)"] \ar[r] & j(Z) \ar[d, "w"] \ar[r, "j(f)"] & j(X)\ar[d] \\
\wt{j}(X)\ar[r, "v"]   & Y \ar[r] & Y'
\end{tikzcd}
\end{equation}
Now the morphism $w:j(Z)\to Y$ is represented by some morphism $g:Z'\to Y,$ where $Z'\subset Z$ and $Z',Z/Z'\in\cB.$ Using the commutativity of \eqref{eq:two_comm_squares} and replacing $Z'$ by a smaller subobject if necessary (so that still $Z/Z'\in\cB$), we may assume that the following square commutes:
\begin{equation*}
\begin{CD}
i(Z') @>{i(g)}>> i(Y)\\
@VVV @V{u}VV\\
Z @>{f}>> X.
\end{CD}
\end{equation*}
By construction, we obtain a well-defined morphism $(f,g):\alpha(Z,Z')\to (X,Y,u,v).$ We also have a natural morphism $\alpha(Y,Y)\to (X,Y,u,v),$ and together these two morphisms define an epimorphism
\begin{equation*}
\alpha(Z,Z')\oplus\alpha(Y,Y)\to (X,Y,u,v).
\end{equation*}
This proves that the object $(X,Y,u,v)$ is a quotient of an object of $\alpha(\cE_{\cA,\cB}),$ as required. 
\end{proof}

Next, we observe that the category $\cA$ is naturally a Serre quotient of $\cC_{\cA,\cB}.$

\begin{prop}\label{prop:Serre_quotient}
The full subcategory $\beta(\cB)\subset \cC_{\cA,\cB}$ is a Serre subcategory, i.e. it is closed under taking subobjects, quotients and extensions. The functor $\pi:\cC_{\cA,\cB}\to \cA,$ $\pi(X,Y,u,v)=X,$ induces an equivalence $\bbar{\pi}:\cC_{\cA,\cB}/\beta(\cB)\xto{\sim} \cA.$
\end{prop}

\begin{proof}
The first statement is evident. Next, we see that the functor $\pi$ is essentially surjective: for any $X\in\cA$ we can choose $Y\subset X$ such that $Y,X/Y\in\cB,$ and then $X\cong \pi(\alpha(X,Y)).$

It remains to prove the fully faithfulness of $\bbar{\pi}.$ By Proposition \ref{prop:properties_of_C_AB} \ref{torsion_pair}, each object of $\cC_{\cA,\cB}/\beta(\cB)$ is isomorphic to an image of an object of $\alpha(\cE_{\cA,\cB}).$ Take some objects $(X_1,Y_1),(X_2,Y_2)\in\cE_{\cA,\cB}.$ If $Z\subset\alpha(X_1,Y_1)$ is a subobject such that $\alpha(X_1,Y_1)/Z\in\beta(\cB),$ then $Z$ must be of the form $\alpha(X_1,Y_1')$ for some $Y_1'\subset Y_1$ such that $X_1/Y_1'\in\cB.$ On the other hand, any non-zero subobject of $\alpha(X_2,Y_2)$ is not contained in $\beta(\cB).$ We obtain
\begin{multline*}
\Hom_{\cC_{\cA,\cB}/\beta(\cB)}(\alpha(X_1,Y_1),\alpha(X_2,Y_2))=\indlim[\substack{Y_1'\subset Y_1,\\ X_1/Y_1'\in\cB}]\Hom_{\cC_{\cA,\cB}}(\alpha(X_1,Y_1'),\alpha(X_2,Y_2))\\
\xto{\sim} \Hom_{\cA}(X_1,X_2).
\end{multline*}
Here the second isomorphism follows from the following: for any $f:X_1\to X_2$ there exists $Y_1'\subset Y_1$ such that $X_1/Y_1'\in\cB$ and $f(Y_1')\subset Y_2.$ Namely, by Lemma \ref{lem:intersections} we can take $Y_1' = Y_1\cap f^{-1}(Y_2).$ This proves the proposition. 
\end{proof}

\section{Proofs of Theorem \ref{th:key_construction_intro} and \ref{th:Quillen}}

\begin{proof}[Proof of Theorem \ref{th:key_construction_intro}]
We prove \ref{SOD}. We first consider the exact functors between quasi-abelian categories
\begin{equation*}
\Phi_1:\cB\to\cE_{\cA,\cB},\,\Phi_1(X)=(X,0),\quad \Phi_2:\cB\to\cE_{\cA,\cB},\,\Phi_2(X)=(X,X).
\end{equation*} 
The left adjoint $\Phi_1^L$ to $\Phi_1$ and the right adjoint $\Phi_2^R$ to $\Phi_2$ are also exact functors, given by
\begin{equation*}
\Phi_1^L(X,Y)=X/Y,\quad \Phi_2^R(X,Y)=Y.
\end{equation*}
Hence, we have a functorial admissible short exact sequence
\begin{equation*}
0\to \Phi_2(\Phi_2^R(X,Y))\to (X,Y)\to \Phi_1(\Phi_1^L(X,Y))\to 0,\quad (X,Y)\in\cE_{\cA,\cB}.
\end{equation*}
We also have the vanishing $\Phi_2^R\circ \Phi_1=0$ and the isomorphisms $\Phi_1^L\circ\Phi_1\xto{\sim}\id,$ $\id\xto{\sim}\Phi_2^R\circ\Phi_2.$ Denoting by the same symbols the induced functors between the bounded derived categories, we see that the functors $\Phi_1:D^b(\cB)\to D^b(\cE_{\cA,\cB})$ and $\Phi_2:D^b(\cB)\to D^b(\cE_{\cA,\cB})$ are fully faithful and we have a semi-orthogonal decomposition
\begin{equation*}
D^b(\cE_{\cA,\cB})=\la \Phi_1(D^b(\cB)),\Phi_2(D^b(\cB))\ra.
\end{equation*}

Next, we prove \ref{ses}. By Proposition \ref{prop:Serre_quotient} and \cite[Lemma 1.15a]{Ke99} we have a short exact sequence of dg categories
\begin{equation}\label{eq:ses_via_C_AB}
0\to D^b_{\beta(\cB)}(\cC_{\cA,\cB})\to D^b(\cC_{\cA,\cB})\xto{\pi} D^b(\cA)\to 0. 
\end{equation}
Here the first category is the full subcategory of complexes whose homology is contained in the essential image of $\beta.$ By Proposition \ref{prop:abelian_envelope} we have an equivalence $D^b(\cE_{\cA,\cB})\simeq D^b(\cC_{\cA,\cB}).$ It induces an equivalence from \eqref{eq:key_ses} to \eqref{eq:ses_via_C_AB}. The assertion about the bounded $t$-structure on $\cT_{\cA,\cB}$ is clear: the composition $\cB\xto{\beta} D^b_{\beta(B)}(\cC_{\cA,\cB})\xto{\sim} \cT_{\cA,\cB}$ is given by
\begin{equation*}
X\mapsto \Cone((X,0)\to (X,X)).
\end{equation*}

Now \ref{functors_on_components} is straightforward. Namely, the isomorphisms $q\circ \Phi_1\cong q\circ \Phi_2\cong i$ follow by construction. The restriction of the composition $\Phi_1^L[-1]\circ\Psi$ to the heart $\cT_{\cA,\cB}^{\heartsuit}\simeq \cB$ is given by the natural inclusion $\cB\to D^b(\cB),$ and similarly for $\Phi_2^R\circ\Psi.$ It follows that both $\Phi_1^L[-1]\circ\Psi$ and $\Phi_2^R\circ\Psi$ are left inverses to the realization functor, as required. 
\end{proof}

In the proof of Theorem \ref{th:Quillen} below we avoid using the vanishing of $K_{-1}(\cT_{\cA,\cB})$ (which is known by Theorem \ref{th:key_construction_intro} \ref{ses} and \cite[Theorem 2.35]{AGH19}), because the statement of Theorem \ref{th:Quillen} for $n=0$ is elementary.

\begin{proof}[Proof of Theorem \ref{th:Quillen}] Let $\cA\supset\cB$ be as in Theorem \ref{th:Quillen}. By a simple argument we reduce to the case when each object of $\cA$ is an extension of two objects of $\cB.$ Indeed, for each $n\geq 1$ let $\cA_n\subset\cA$ be the full subcategory which consists of objects with a finite filtration with at most $n$ non-zero subquotients which are contained in $\cB.$ Then each $\cA_n$ is closed under taking subobjects and quotients, $\cA_1=\cB$ and $\bigcup\limits_{n}\cA_n=\cA.$ Since $K$-theory commutes with filtered colimits, we may assume that $\cA=\cA_n$ for some $n.$ Using induction we reduce to the case $\cA=\cA_2.$ We assume this from now on.

To show the isomorphism $K_0(\cB)\xto{\sim} K_0(\cA),$ we define the inverse map $\gamma:K_0(\cA)\to K_0(\cB)$ by putting $\gamma([X])=[Y]+[X/Y]$ for some $Y\subset X$ such that $Y,X/Y\in \cB.$ We only need to check that $\gamma$ is well-defined. To see this, first take another subobject $Y'\subset X$ such that $Y',X/Y'\in\cB.$ Then we have
\begin{multline*}
[Y]+[X/Y]=[Y\cap Y']+[Y/(Y\cap Y')]+[X/Y]=[Y\cap Y']+[X/(Y\cap Y')]\\
=[Y\cap Y']+[Y'/(Y\cap Y')]+[X/Y']=[Y']+[X/Y'].
\end{multline*} 
For a short exact sequence 
\begin{equation*}
0\to X'\xto{f} X\xto{g} X''\to 0
\end{equation*}
in $\cA,$ take some $Y\subset X$ such that $Y,X/Y\in\cB,$ and put $Y'=f^{-1}(Y)\subset X'$ and $Y''=g(Y)\subset X''.$ Then we have
\begin{equation*}
\gamma([X])-\gamma([X'])-\gamma([X''])=([Y]-[Y']+[Y''])+([X/Y]-[X'/Y']-[X''/Y''])=0.
\end{equation*}
Hence, $\gamma$ is well-defined and $K_0(\cB)\cong K_0(\cA).$

To prove the isomorphisms $K_n(\cB)\xto{\sim}K_n(\cA)$ for $n\geq 1$ we use the localization theorem for the short exact sequence \eqref{eq:key_ses}. By Theorem \ref{th:Barwick} the realization functor induces isomorphisms $K_n(\cB)=K_n(D^b(\cB))\xto{\sim} K_n(\cT_{\cA,\cB})$ for $n\geq 0.$ By Theorem \ref{th:key_construction_intro} \ref{functors_on_components}, for $n\geq 0$ the composition
\begin{equation*}
K_n(\cB)\cong K_n(\cT_{\cA,\cB})\to K_n(D^b(\cE_{\cA,\cB}))\cong K_n(\cB)\oplus K_n(\cB)
\end{equation*}
is given by $(-\id,\id).$ Therefore, the long exact sequence of $K$-groups for \eqref{eq:key_ses} implies the isomorphisms $K_n(\cB)\xto{\sim} K_n(\cA)$ for $n\geq 1.$
\end{proof}

\begin{thebibliography}{BVdB03}

%\bibitem[AR94]{AR}J.~Ad\'amek, J. Rosick\'y. ``Locally presentable and accessible categories''. London Math. Society
%Lecture Note Series 189, Cambridge University Press, 1994

\bibitem[AGH19]{AGH19} B.~Antieau, D.~Gepner and J.~Heller, ``$K$-theoretic obstructions to bounded $t$-structures''. Invent. math. 216, 241-300 (2019). 

\bibitem[Bar15]{Bar15} C.~Barwick, ``On exact $\infty$-categories and the Theorem of the Heart''. Compos. Math. 151(11), 2160–2186 (2015).

%\bibitem[Bar16]{Bar}C.~Barwick, ``On the algebraic K-theory of higher categories''. J. Topol. 9 (2016), no. 1,
%245-347.

\bibitem[BVdB03]{BVdB03} A.~Bondal and M.~Van den Bergh, ``Generators and representability of functors in commutative and
noncommutative geometry''. Mosc. Math. J. 3 (2003), no. 1, 1-36, 258.

\bibitem[E20]{E20} A.~I.~Efimov, ``Homotopy finiteness of some DG categories from algebraic geometry''. J. Eur. Math. Soc. 22 (2020), no. 9, pp. 2879-2942.

%\bibitem[E24]{E24}A.~I.~Efimov, ``$K$-theory and localizing invariants of large categories'',  arXiv:2405.12169 (preprint).

\bibitem[GP64]{GP64} P.~Gabriel and N.~Popesco, ``Caract\'erisation des cat\'egories ab\'eliennes avec g\'en\'erateurs et
limites inductives exactes''. C. R. Acad. Sci. Paris 258 (1964), 4188-4190.

\bibitem[HRS96]{HRS96} D.~Happel, I.~Reiten, and S.~Smalo, ``Tilting in abelian categories and quasitilted algebra''. Memoirs of
the AMS, vol. 575, Amer. Math. Soc., 1996.

\bibitem[Ka22]{Ka22} D.~B.~Kaledin, ``What do abelian categories form?''. Russian Math. Surveys, 77:1 (2022), 1–45. 

\bibitem[KS01]{KS01} M.~Kashiwara, S.~Schapira, ``Ind-sheaves'', Ast\'erisque, vol. 271, Soc. Math. France, 2001.

\bibitem[Ke99]{Ke99} B.~Keller, ``On the cyclic homology category of exact categories''. J. Pure Appl. Algebra 136, 1–56 (1999).

\bibitem[Ku94]{Ku94} N.~J.~Kuhn, ``Generic representations of the finite general linear groups and the Steenrod algebra. I''. Amer. J. Math. 116 (1994), no. 2, 327-360.

\bibitem[KL15]{KL15} A.~Kuznetsov and V.~Lunts, ``Categorical resolutions of irrational singularities''. Int. Math. Res. Notices 2015, 4536-4625.

\bibitem[LT19]{LT19} M.~Land and G.~Tamme, ``On the K-theory of pullbacks''. Ann. of Math. (2) 190 (2019), no. 3, 877–930. 

\bibitem[Sch99]{Sch99} J.-P. Schneiders, ``Quasi-abelian categories and sheaves''. M\'em. Soc. Math. Fr. (N.S.) (1999), no. 76,
vi+134.

\bibitem[Qui73]{Qui} D.~Quillen, ``Higher algebraic K-theory. I''. Algebraic K-theory, I: Higher K-theories
(Proc. Conf., Battelle Memorial Inst., Seattle, Wash., 1972), Springer, Berlin, 1973,
pp. 85-147. Lecture Notes in Math., Vol. 341.


\end{thebibliography}
\end{document}
